\section{Proofs for \texorpdfstring{Section~\ref{sec:evo}}{evolutionary individuality}}\label{app:evo}

The statements of \cref{sec:evo} are proved in their order of appearance, after two lemmas on $2\times2$ matrices. Three secondary results (\cref{prop:evo-bottleneck,prop:evo-directional,prop:evo-two-type}), a worked example and an extended comparison with \citet{bourrat2022tradeoff} follow. Throughout, $\tilde\varepsilon=\varepsilon+g(1-\mathsf f)$ and $\tilde\beta=\beta+g\mathsf f$ as in \cref{cor:hie-identity}, and $\mathsf p_1=\varepsilon+\beta+g=\tilde\varepsilon+\tilde\beta$, $\mathsf p_2=\omega+\delta$, $\mathsf E=\omega\tilde\beta$ and $\varpi_\pm$ are as in the proof of \cref{prop:agy-identity}, so that $\varpi_-\le\varpi_+$ are the roots of
\begin{equation}\label{eq:evo-charpoly}
\chi(s)=(\mathsf p_1-s)(\mathsf p_2-s)-\mathsf E,
\end{equation}
and $\tau_{\mathrm{id}}^{-1}=\varpi_-$. The standing hypotheses are those of \cref{sec:evo}: the hypotheses of \cref{prop:agy-identity}, among them $g>0$, $\omega>0$ and $\det\mathsf A=\mathsf p_1\mathsf p_2-\mathsf E=\omega\tilde\varepsilon+\delta\mathsf p_1>0$, together with $\tilde\beta>0$ and $\mathsf T>0$. Of these, \cref{lem:evo-semigroup} uses only $\omega>0$, $\tilde\beta>0$ and $\det\mathsf A>0$. Inequalities between vectors or between matrices are entrywise, $\one=(1,1)^\top$, and $\mathsf G(t)=e^{t\mathsf A}$ with entries $\mathsf G_{ij}(t)$.

\begin{lemma}[Nonnegative $2\times2$ matrices]\label{lem:evo-perron}
Let $\mathsf Q=(q_{ij})$ be a real $2\times2$ matrix with nonnegative entries and put $\lambda_\pm=\frac12[q_{11}+q_{22}\pm((q_{11}-q_{22})^2+4q_{12}q_{21})^{1/2}]$.
\begin{enumerate}
\item[(i)] The eigenvalues of $\mathsf Q$ are $\lambda_\pm$; they are real, $\lambda_+\ge\max\{q_{11},q_{22}\}$ and $\lambda_+\ge\lvert\lambda_-\rvert$. Hence the spectral radius is $\rho(\mathsf Q)=\lambda_+$, a continuous function of the entries.
\item[(ii)] There are nonzero vectors $w\ge0$ and $v\ge0$ with $\mathsf Qw=\lambda_+w$ and $v^\top\mathsf Q=\lambda_+v^\top$.
\item[(iii)] If $\vartheta>0$ and $a\vartheta^\top\le\vartheta^\top\mathsf Q\le b\vartheta^\top$ for real numbers $a\le b$, then $a\le\lambda_+\le b$. If $\mathsf Q\one\le\one$ then $\lambda_+\le1$, and if $\mathsf Q\one<\one$ then $\lambda_+<1$.
\item[(iv)] If all entries are positive, then $\lambda_+>\lvert\lambda_-\rvert$, the vectors $w$ and $v$ of (ii) are positive and unique up to scale, and $\lambda_+^{-k}\mathsf Q^k\to wv^\top/(v^\top w)$ as $k\to\infty$.
\end{enumerate}
\end{lemma}

\begin{proof}
(i) The characteristic polynomial is $\lambda^2-(q_{11}+q_{22})\lambda+q_{11}q_{22}-q_{12}q_{21}$, with discriminant $D=(q_{11}-q_{22})^2+4q_{12}q_{21}\ge(q_{11}-q_{22})^2\ge0$. Its roots are $\lambda_\pm$, they are real, and $\lambda_+\ge\frac12[q_{11}+q_{22}+\lvert q_{11}-q_{22}\rvert]=\max\{q_{11},q_{22}\}$. Since $\lambda_++\lambda_-=q_{11}+q_{22}\ge0$ and $\lambda_+\ge\lambda_-$, also $\lambda_+\ge-\lambda_-$, so $\lambda_+\ge\lvert\lambda_-\rvert$. Continuity is read off the formula.

(ii) The characteristic equation at $\lambda_+$ reads $(\lambda_+-q_{11})(\lambda_+-q_{22})=q_{12}q_{21}$. If $q_{12}>0$, the vector $w=(q_{12},\lambda_+-q_{11})^\top$ is nonzero and nonnegative by (i), the first row of $\mathsf Qw$ is $q_{11}q_{12}+q_{12}(\lambda_+-q_{11})=\lambda_+q_{12}$, and the second row is $q_{21}q_{12}+q_{22}(\lambda_+-q_{11})=(\lambda_+-q_{11})(\lambda_+-q_{22})+q_{22}(\lambda_+-q_{11})=\lambda_+(\lambda_+-q_{11})$. If $q_{12}=0$ and $q_{22}\ge q_{11}$, then $\lambda_+=q_{22}$ and $w=(0,1)^\top$ works. If $q_{12}=0$ and $q_{11}>q_{22}$, then $\lambda_+=q_{11}$ and $w=(q_{11}-q_{22},q_{21})^\top$ works, because the second row of $\mathsf Qw$ is $q_{21}(q_{11}-q_{22})+q_{22}q_{21}=q_{11}q_{21}$. The vector $v$ is obtained by applying this to $\mathsf Q^\top$, which has the same $\lambda_+$.

(iii) With $w$ from (ii), $\vartheta^\top w>0$, and multiplying the hypothesis on the right by $w\ge0$ gives $a\vartheta^\top w\le\vartheta^\top\mathsf Qw=\lambda_+\vartheta^\top w\le b\vartheta^\top w$. With $v$ from (ii), $v^\top\one>0$ and $\lambda_+v^\top\one=v^\top\mathsf Q\one\le v^\top\one$, with strict inequality when $\mathsf Q\one<\one$ because $v\ge0$ is nonzero.

(iv) Now $D>0$, so $\lambda_+>\lambda_-$, and $\lambda_++\lambda_->0$, so $\lambda_+>-\lambda_-$. In the characteristic equation the right side is positive and both factors on the left are nonnegative by (i), so both are positive; hence $w=(q_{12},\lambda_+-q_{11})^\top$ and $v=(q_{21},\lambda_+-q_{11})^\top$ are positive. The eigenvalues are distinct, so each eigenvector is unique up to scale and $\mathsf Q$ is diagonalizable. If $w_-$ is a right eigenvector of $\lambda_-$, then $\lambda_+v^\top w_-=v^\top\mathsf Qw_-=\lambda_-v^\top w_-$ gives $v^\top w_-=0$, so $\Pi=wv^\top/(v^\top w)$ fixes $w$ and annihilates $w_-$, and $\mathsf Q^k=\lambda_+^k\Pi+\lambda_-^k(I-\Pi)$. Dividing by $\lambda_+^k$ and using $\lvert\lambda_-/\lambda_+\rvert<1$ gives the limit.
\end{proof}

\begin{lemma}[The relaxation semigroup]\label{lem:evo-semigroup}
Under the standing hypotheses the following hold.
\begin{enumerate}
\item[(i)] $0<\varpi_-<\min\{\mathsf p_1,\mathsf p_2\}\le\max\{\mathsf p_1,\mathsf p_2\}<\varpi_+$, and
\begin{equation}\label{eq:evo-root-identities}
\varpi_+-\mathsf p_2=\mathsf p_1-\varpi_-,\qquad\varpi_+-\mathsf p_1=\mathsf p_2-\varpi_-,\qquad(\mathsf p_1-\varpi_-)(\mathsf p_2-\varpi_-)=\mathsf E.
\end{equation}
The number $\mathsf u=(\mathsf p_1-\varpi_-)/(\varpi_+-\varpi_-)=\frac12[1+(\mathsf p_1-\mathsf p_2)((\mathsf p_1-\mathsf p_2)^2+4\mathsf E)^{-1/2}]$ lies in $(0,1)$, and $1-\mathsf u=(\mathsf p_2-\varpi_-)/(\varpi_+-\varpi_-)$.
\item[(ii)] For every $t\ge0$,
\begin{equation}\label{eq:evo-semigroup}
\begin{aligned}
\mathsf G(t)&=\frac{e^{-\varpi_-t}(\mathsf A+\varpi_+I)-e^{-\varpi_+t}(\mathsf A+\varpi_-I)}{\varpi_+-\varpi_-}\\
&=\begin{pmatrix}(1-\mathsf u)e^{-\varpi_-t}+\mathsf ue^{-\varpi_+t}&\tilde\beta\dfrac{e^{-\varpi_-t}-e^{-\varpi_+t}}{\varpi_+-\varpi_-}\\[2ex]\omega\dfrac{e^{-\varpi_-t}-e^{-\varpi_+t}}{\varpi_+-\varpi_-}&\mathsf ue^{-\varpi_-t}+(1-\mathsf u)e^{-\varpi_+t}\end{pmatrix},
\end{aligned}
\end{equation}
and all entries of $\mathsf G(t)$ are positive for $t>0$.
\item[(iii)] The vectors $\ell=(\omega,\mathsf p_1-\varpi_-)^\top$ and $r=(\tilde\beta,\mathsf p_1-\varpi_-)^\top$ are positive, $\ell^\top\mathsf A=-\varpi_-\ell^\top$, $\mathsf Ar=-\varpi_-r$, $\ell^\top r=(\mathsf p_1-\varpi_-)(\varpi_+-\varpi_-)$, and $e^{\varpi_-t}\mathsf G(t)\to r\ell^\top/(\ell^\top r)$ as $t\to\infty$.
\item[(iv)] $\mathsf G(t)\one<\one$ for every $t>0$.
\item[(v)] $\chi(\tilde\varepsilon)=\tilde\beta(\delta-\tilde\varepsilon)$ and $\chi(\delta)=\omega(\tilde\varepsilon-\delta)$. Consequently $\min\{\delta,\tilde\varepsilon\}\le\varpi_-\le\max\{\delta,\tilde\varepsilon\}$, both inequalities are strict when $\delta\ne\tilde\varepsilon$, and $\varpi_-=\delta=\tilde\varepsilon$ when $\delta=\tilde\varepsilon$.
\end{enumerate}
\end{lemma}

\begin{proof}
(i) The polynomial $\chi$ of \eqref{eq:evo-charpoly} is a monic quadratic with real roots $\varpi_-\le\varpi_+$ and is negative exactly on $(\varpi_-,\varpi_+)$. Since $\chi(\mathsf p_1)=\chi(\mathsf p_2)=-\mathsf E<0$, both $\mathsf p_1$ and $\mathsf p_2$ lie strictly between the roots. The product of the roots is $\det\mathsf A>0$ and $\varpi_+>0$, so $\varpi_->0$. The sum of the roots is $\mathsf p_1+\mathsf p_2$, which gives the first two identities, and $\chi(\varpi_-)=0$ is the third. Then $\mathsf u\in(0,1)$ because $\varpi_-<\mathsf p_1<\varpi_+$, the expression for $1-\mathsf u$ follows from the second identity, and the closed form of $\mathsf u$ from $\varpi_+-\varpi_-=((\mathsf p_1-\mathsf p_2)^2+4\mathsf E)^{1/2}$ and $\mathsf p_1-\varpi_-=\frac12(\mathsf p_1-\mathsf p_2)+\frac12(\varpi_+-\varpi_-)$.

(ii) Let $\mathsf X(t)$ be the middle expression of \eqref{eq:evo-semigroup}. Then $\mathsf X(0)=I$. By the Cayley--Hamilton theorem $(\mathsf A+\varpi_-I)(\mathsf A+\varpi_+I)=0$, so $\mathsf A(\mathsf A+\varpi_+I)=-\varpi_-(\mathsf A+\varpi_+I)$ and $\mathsf A(\mathsf A+\varpi_-I)=-\varpi_+(\mathsf A+\varpi_-I)$, whence $\mathsf A\mathsf X(t)=\mathrm d\mathsf X(t)/\mathrm dt$. The linear equation $\mathrm d\mathsf X/\mathrm dt=\mathsf A\mathsf X$ with $\mathsf X(0)=I$ has the unique solution $e^{t\mathsf A}$. The entries follow from
\begin{equation*}
\mathsf A+\varpi_+I=\begin{pmatrix}\mathsf p_2-\varpi_-&\tilde\beta\\\omega&\mathsf p_1-\varpi_-\end{pmatrix},\qquad\mathsf A+\varpi_-I=\begin{pmatrix}-(\mathsf p_1-\varpi_-)&\tilde\beta\\\omega&-(\mathsf p_2-\varpi_-)\end{pmatrix},
\end{equation*}
where \eqref{eq:evo-root-identities} was used in the first matrix. For $t>0$ one has $e^{-\varpi_-t}>e^{-\varpi_+t}>0$, and $\omega$, $\tilde\beta$, $\mathsf u$ and $1-\mathsf u$ are positive, so every entry is positive.

(iii) Positivity is (i). The first component of $\ell^\top\mathsf A$ is $-\mathsf p_1\omega+\omega(\mathsf p_1-\varpi_-)=-\varpi_-\omega$, and the second is $\tilde\beta\omega-\mathsf p_2(\mathsf p_1-\varpi_-)=(\mathsf p_1-\varpi_-)(\mathsf p_2-\varpi_-)-\mathsf p_2(\mathsf p_1-\varpi_-)=-\varpi_-(\mathsf p_1-\varpi_-)$ by \eqref{eq:evo-root-identities}; the computation for $r$ is the same with $\omega$ and $\tilde\beta$ exchanged. Next, $\ell^\top r=\mathsf E+(\mathsf p_1-\varpi_-)^2=(\mathsf p_1-\varpi_-)(\mathsf p_1+\mathsf p_2-2\varpi_-)=(\mathsf p_1-\varpi_-)(\varpi_+-\varpi_-)$. Finally, using $\omega\tilde\beta=(\mathsf p_1-\varpi_-)(\mathsf p_2-\varpi_-)$, one checks entry by entry that $r\ell^\top=(\mathsf p_1-\varpi_-)(\mathsf A+\varpi_+I)$, so by \eqref{eq:evo-semigroup} $e^{\varpi_-t}\mathsf G(t)=r\ell^\top/(\ell^\top r)-e^{-(\varpi_+-\varpi_-)t}(\mathsf A+\varpi_-I)/(\varpi_+-\varpi_-)$, and the second term vanishes as $t\to\infty$.

(iv) $\mathsf A\one=(-\mathsf p_1+\tilde\beta,\omega-\mathsf p_2)^\top=(-\tilde\varepsilon,-\delta)^\top\le0$, and this vector is nonzero because $\det\mathsf A=\omega\tilde\varepsilon+\delta\mathsf p_1>0$. Since $\mathsf G(s)>0$ for $s>0$ by (ii), $\mathrm d(\mathsf G(s)\one)/\mathrm ds=\mathsf G(s)\mathsf A\one<0$ for $s>0$, and integrating from $0$ to $t$ gives $\mathsf G(t)\one<\one$.

(v) $\chi(\tilde\varepsilon)=\tilde\beta(\omega+\delta-\tilde\varepsilon)-\omega\tilde\beta=\tilde\beta(\delta-\tilde\varepsilon)$ and $\chi(\delta)=(\tilde\varepsilon+\tilde\beta-\delta)\omega-\omega\tilde\beta=\omega(\tilde\varepsilon-\delta)$. By (i), $\tilde\varepsilon<\mathsf p_1<\varpi_+$ and $\delta<\mathsf p_2<\varpi_+$, so for $s\in\{\tilde\varepsilon,\delta\}$ the sign of $\chi(s)$ decides the position of $s$ relative to $\varpi_-$: $s<\varpi_-$ when $\chi(s)>0$, $s>\varpi_-$ when $\chi(s)<0$, and $s=\varpi_-$ when $\chi(s)=0$. If $\delta<\tilde\varepsilon$ then $\chi(\tilde\varepsilon)<0<\chi(\delta)$, so $\delta<\varpi_-<\tilde\varepsilon$. If $\delta>\tilde\varepsilon$ the signs are reversed and $\tilde\varepsilon<\varpi_-<\delta$. If $\delta=\tilde\varepsilon$ both values vanish.
\end{proof}

\begin{proof}[Proof of \cref{prop:evo-transmission}]
Let $\mathcal H$ be the parental information of \cref{def:evo-fission}, and for one daughter let $\xi_i\in\{0,1\}$ indicate that position $i$ is a founder.

(a) The daughter's mean at birth is $\bar\mu'_0=\frac1n\sum_i[\xi_i\mu_i+(1-\xi_i)(\mathsf f z'_0+\nu_i)]$ and its record is $z'_0=\mathsf cz+\zeta$. The founder set is independent of $\mathcal H$ with $\E\xi_i=\bar{\mathsf b}$, and $\zeta$ and the $\nu_i$ have zero mean and are independent of $\mathcal H$ and of the founder set, so $\E[z'_0\mid\mathcal H]=\mathsf cz$ and $\E[\bar\mu'_0\mid\mathcal H]=\bar{\mathsf b}\bar\mu+(1-\bar{\mathsf b})\mathsf f\mathsf cz$; that is, $\E[\mathsf x'_0\mid\mathcal H]=\mathsf F\mathsf x$. By step (iv) of \cref{def:evo-fission} the daughter's noise increments after birth have zero conditional mean given $\mathcal H$ and the daughter's own past, so the argument of \cref{prop:agy-identity}(a) applies to the daughter with $\mathcal H$ adjoined to the conditioning set and gives $\E[\mathsf x'(\mathsf T)\mid\mathcal H,\mathsf x'_0]=\mathsf G(\mathsf T)\mathsf x'_0$. The tower property gives $\E[\mathsf x'(\mathsf T)\mid\mathcal H]=\mathsf G(\mathsf T)\mathsf F\mathsf x=\mathsf M\mathsf x$. For a parent drawn from a population with finite second moments and covariance $\Sigma$ of $\mathsf x$, the parent's state is $\mathcal H$-measurable, so $\Cov(\mathsf x'(\mathsf T),\mathsf x)=\Cov(\E[\mathsf x'(\mathsf T)\mid\mathcal H],\mathsf x)=\mathsf M\Sigma$, and the regression of $\mathsf a^\top\mathsf x'(\mathsf T)$ on $\mathsf a^\top\mathsf x$ has slope $\mathsf a^\top\mathsf M\Sigma\mathsf a/(\mathsf a^\top\Sigma\mathsf a)$ whenever the denominator is positive.

(b) By \cref{lem:evo-semigroup}(ii) $\mathsf G(\mathsf T)>0$, and $\mathsf F\ge0$, so $\mathsf M\ge0$, with entries
\begin{equation}\label{eq:evo-M-entries}
\mathsf M_{11}=\bar{\mathsf b}\mathsf G_{11},\quad\mathsf M_{21}=\bar{\mathsf b}\mathsf G_{21},\quad\mathsf M_{12}=\mathsf c[(1-\bar{\mathsf b})\mathsf f\mathsf G_{11}+\mathsf G_{12}],\quad\mathsf M_{22}=\mathsf c[(1-\bar{\mathsf b})\mathsf f\mathsf G_{21}+\mathsf G_{22}],
\end{equation}
all evaluated at $\mathsf T$. \Cref{lem:evo-perron}(i) and (ii) give the formula \eqref{eq:evo-heritability} for $\mathsf h=\rho(\mathsf M)$ and a nonzero $\mathsf l\ge0$ with $\mathsf l^\top\mathsf M=\mathsf h\mathsf l^\top$. Then $\E[y'\mid\mathcal H]=\mathsf l^\top\mathsf M\mathsf x=\mathsf hy$ by (a), and the slope of (a) with $\mathsf a=\mathsf l$ is $\mathsf h\mathsf l^\top\Sigma\mathsf l/(\mathsf l^\top\Sigma\mathsf l)=\mathsf h$. Since $\mathsf F\one=(\bar{\mathsf b}+(1-\bar{\mathsf b})\mathsf f\mathsf c,\mathsf c)^\top\le\one$ and $\mathsf G(\mathsf T)\ge0$, \cref{lem:evo-semigroup}(iv) gives $\mathsf M\one\le\mathsf G(\mathsf T)\one<\one$, and \cref{lem:evo-perron}(iii) gives $\mathsf h<1$. When $\mathsf b=0$ the first column of $\mathsf M$ vanishes, $\mathsf h=\mathsf M_{22}$ and $\mathsf l=(0,1)^\top$; when $\mathsf c=0$ the second column vanishes, $\mathsf h=\mathsf M_{11}=\bar{\mathsf b}\mathsf G_{11}(\mathsf T)$ and $\mathsf l=(1,0)^\top$ if $\mathsf b\ge1$.

(c) If $\mathsf b\ge1$ and $\mathsf c>0$, every entry in \eqref{eq:evo-M-entries} is positive, so \cref{lem:evo-perron}(iv) applies: $\mathsf l>0$ is unique up to scale and $\mathsf h^{-k}\mathsf M^k\to r_{\mathsf M}\mathsf l^\top/(\mathsf l^\top r_{\mathsf M})$, where $r_{\mathsf M}>0$ is the right eigenvector of $\mathsf h$. Consider a line of descent in which every collective founds at least one daughter, so that the line exists for every $k$, and in which the daughter followed is selected by a rule independent of the states and noises of all collectives on the line. Let $\mathcal H_k$ be the information of generation $k$ on the line, with $\mathcal H_0=\mathcal H$ and $\mathcal H_k\subseteq\mathcal H_{k+1}$. Because the selection rule is independent of the daughters' states and noises, conditioning on which daughter is followed does not change the conditional law of that daughter's state, and (a) applied to generation $k$ gives $\E[\mathsf x^{(k+1)}\mid\mathcal H_k]=\mathsf M\mathsf x^{(k)}$. The tower property and induction on $k$ give $\E[\mathsf x^{(k)}\mid\mathcal H]=\mathsf M^k\mathsf x$, so $\mathsf h^{-k}\E[\mathsf a^\top\mathsf x^{(k)}\mid\mathcal H]\to(\mathsf a^\top r_{\mathsf M})(\mathsf l^\top\mathsf x)/(\mathsf l^\top r_{\mathsf M})$, which is the stated limit; it is nonzero exactly when $\mathsf a^\top r_{\mathsf M}\ne0$ and $y\ne0$, since $\mathsf l^\top r_{\mathsf M}>0$. A rule that selects the daughter by its state, for example the daughter with the largest trait, is excluded: the selected daughter's conditional mean is then not $\mathsf M\mathsf x^{(k)}$ in general. If $\mathsf a^\top r_{\mathsf M}=0$ with $\mathsf a\ne0$, then $\mathsf a$ is a left eigenvector of $\mathsf M$ for its second eigenvalue $\lambda_-=\tr\mathsf M-\mathsf h$, because the eigenvalues are distinct and the left eigenvector of one annihilates the right eigenvector of the other, and $\E[\mathsf a^\top\mathsf x^{(k)}\mid\mathcal H]=\lambda_-^k\mathsf a^\top\mathsf x$ with $\lvert\lambda_-\rvert<\mathsf h$.

(d) In the variant without a record, $\bar\mu'_0=\frac1n\sum_i[\xi_i\mu_i+(1-\xi_i)(\mathsf f m_S+\nu_i)]=(\bar{\mathsf b}+(1-\bar{\mathsf b})\mathsf f)m_S+\frac1n\sum_i(1-\xi_i)\nu_i$, where $m_S=\frac1{\mathsf b}\sum_i\xi_i\mu_i$ has conditional mean $\bar\mu$ given $\mathcal H$. By the last sentence of \cref{prop:agy-identity}(a) the conditional mean of the members' mean then relaxes at the rate $\tilde\varepsilon$, so $\E[\bar\mu'(\mathsf T)\mid\mathcal H]=(\bar{\mathsf b}+(1-\bar{\mathsf b})\mathsf f)e^{-\tilde\varepsilon\mathsf T}\bar\mu=\mathsf h_0\bar\mu$.
\end{proof}

\begin{proof}[Proof of \cref{thm:evo-record-channel}]
(a) With $\ell>0$ from \cref{lem:evo-semigroup}(iii), $\ell^\top\mathsf M=\ell^\top\mathsf G(\mathsf T)\mathsf F=e^{-\varpi_-\mathsf T}\ell^\top\mathsf F$, and
\begin{equation}\label{eq:evo-lF}
\ell^\top\mathsf F=\bigl(\bar{\mathsf b}\omega,\ \omega(1-\bar{\mathsf b})\mathsf f\mathsf c+\mathsf c(\mathsf p_1-\varpi_-)\bigr)=(\bar{\mathsf b}\ell_1,\mathsf c\mathsf k\ell_2),\qquad\mathsf k=1+\frac{(1-\bar{\mathsf b})\mathsf f\omega}{\mathsf p_1-\varpi_-}.
\end{equation}
Hence $e^{-\varpi_-\mathsf T}\min\{\bar{\mathsf b},\mathsf c\mathsf k\}\ell^\top\le\ell^\top\mathsf M\le e^{-\varpi_-\mathsf T}\max\{\bar{\mathsf b},\mathsf c\mathsf k\}\ell^\top$, and \cref{lem:evo-perron}(iii) gives the two-sided bound of \eqref{eq:evo-sandwich}. By \cref{lem:evo-perron}(i), $\mathsf h\ge\max\{\mathsf M_{11},\mathsf M_{22}\}$, and by \eqref{eq:evo-M-entries} and \eqref{eq:evo-semigroup}, $\mathsf M_{11}\ge\bar{\mathsf b}(1-\mathsf u)e^{-\varpi_-\mathsf T}$ and $\mathsf M_{22}\ge\mathsf c\mathsf G_{22}(\mathsf T)\ge\mathsf c\mathsf ue^{-\varpi_-\mathsf T}$, which is the third bound. For the limits, $e^{\varpi_-\mathsf T}\mathsf M\to r\ell^\top\mathsf F/(\ell^\top r)$ as $\mathsf T\to\infty$ by \cref{lem:evo-semigroup}(iii). The limit is a nonnegative matrix of rank at most one, whose spectral radius is its trace $\ell^\top\mathsf Fr/(\ell^\top r)$, and the spectral radius is continuous (\cref{lem:evo-perron}(i)), so by \eqref{eq:evo-lF} and \eqref{eq:evo-root-identities}
\begin{equation*}
\lim_{\mathsf T\to\infty}e^{\varpi_-\mathsf T}\mathsf h=\frac{\bar{\mathsf b}\omega\tilde\beta+\mathsf c\mathsf k(\mathsf p_1-\varpi_-)^2}{(\mathsf p_1-\varpi_-)(\varpi_+-\varpi_-)}=\bar{\mathsf b}\frac{\mathsf p_2-\varpi_-}{\varpi_+-\varpi_-}+\mathsf c\mathsf k\frac{\mathsf p_1-\varpi_-}{\varpi_+-\varpi_-}=(1-\mathsf u)\bar{\mathsf b}+\mathsf u\mathsf c\mathsf k.
\end{equation*}
As $\mathsf T\to0$, $\mathsf M\to\mathsf F$, which is triangular with spectral radius $\max\{\bar{\mathsf b},\mathsf c\}$.

(b) This is \cref{lem:evo-semigroup}(v), with $\tau_{\mathrm{id}}^{-1}=\varpi_-$ and the identity time $1/\tilde\varepsilon$ that \cref{prop:agy-identity}(a) gives without a record.

(c) By (a), $\mathsf h\ge\mathsf c\mathsf ue^{-\varpi_-\mathsf T}>0$ for every $\mathsf b$, while $\mathsf h_0=(\bar{\mathsf b}+(1-\bar{\mathsf b})\mathsf f)e^{-\tilde\varepsilon\mathsf T}$ for $\mathsf b\ge1$ and $\mathsf h_0=0$ for $\mathsf b=0$. Dividing gives the bound, and $\tilde\varepsilon-\varpi_->0$ by (b) when $\delta<\tilde\varepsilon$, so the bound diverges with $\mathsf T$; the weaker form uses $\bar{\mathsf b}+(1-\bar{\mathsf b})\mathsf f\le1$.

(d) By (a), $\mathsf h/\mathsf h_0\le\max\{\bar{\mathsf b},\mathsf c\mathsf k\}e^{-(\varpi_--\tilde\varepsilon)\mathsf T}/(\bar{\mathsf b}+(1-\bar{\mathsf b})\mathsf f)$, and $\varpi_-\ge\tilde\varepsilon$ by (b) when $\delta\ge\tilde\varepsilon$, with strict inequality when $\delta>\tilde\varepsilon$. The prefactor does not depend on $\mathsf T$. For the rate, when $\mathsf b\ge1$ the limit in (a) is at least $(1-\mathsf u)\bar{\mathsf b}>0$, so $\frac1{\mathsf T}\log\mathsf h\to-\varpi_-$, while $\frac1{\mathsf T}\log\mathsf h_0\to-\tilde\varepsilon$. If $\mathsf c=0$, the second column of $\mathsf M$ vanishes and $\mathsf h=\mathsf M_{11}=\bar{\mathsf b}\mathsf G_{11}(\mathsf T)\le\bar{\mathsf b}$, because $\mathsf G_{11}(\mathsf T)$ is a convex combination of $e^{-\varpi_-\mathsf T}$ and $e^{-\varpi_+\mathsf T}$.

(e) Let the record be written only by non-members whose writes do not depend on the members' states, with locally integrable path $z_{\mathrm{ext}}(\cdot)$. As in \cref{prop:agy-identity}(b), the path is then a function of the exterior's exogenous input, which is the conditioning set of \cref{def:agy-identity}; the argument below uses this only through the three conditional hypotheses of the statement. The hypotheses are stated conditionally on the pair $(\mathcal H,z_{\mathrm{ext}})$, since independence of the path from the founder set and from the noises separately would not suffice once the path is added to the conditioning set: given $\mathcal H$ and $z_{\mathrm{ext}}$ the founder set is uniform, so $\E[\xi_i\mid\mathcal H,z_{\mathrm{ext}}]=\bar{\mathsf b}$, the entry noises have $\E[\nu_i\mid\mathcal H,z_{\mathrm{ext}},\xi]=0$, and the relaxation noise of the daughter has zero conditional mean increments given $\mathcal H$, $z_{\mathrm{ext}}$ and the daughter's past. Recruits at fission start at $\mathsf f z_{\mathrm{ext}}(0)+\nu_i$, so $\E[\bar\mu'_0\mid\mathcal H,z_{\mathrm{ext}}]=\bar{\mathsf b}\bar\mu+(1-\bar{\mathsf b})\mathsf f z_{\mathrm{ext}}(0)$. As in the proof of \cref{prop:agy-identity}(b), the first row of \eqref{eq:agy-identity} holds with the record path given: the member-level drifts sum to $-\mathsf p_1\bar\mu'+\tilde\beta z_{\mathrm{ext}}$ because the communication term has zero sum over members, and the noise increments drop out of the conditional mean by the third hypothesis. Hence $m(t)=\E[\bar\mu'(t)\mid\mathcal H,z_{\mathrm{ext}}]$ solves $\mathrm dm/\mathrm dt=-\mathsf p_1m+\tilde\beta z_{\mathrm{ext}}(t)$, and
\begin{equation*}
\E[\bar\mu'(\mathsf T)\mid\mathcal H,z_{\mathrm{ext}}]=\bar{\mathsf b}e^{-\mathsf p_1\mathsf T}\bar\mu+e^{-\mathsf p_1\mathsf T}(1-\bar{\mathsf b})\mathsf f z_{\mathrm{ext}}(0)+\tilde\beta\int_0^{\mathsf T}e^{-\mathsf p_1(\mathsf T-s)}z_{\mathrm{ext}}(s)\mathrm ds.
\end{equation*}
The last two terms are functions of the record path alone, so the coefficient of the parent's mean given the record path is $\bar{\mathsf b}e^{-\mathsf p_1\mathsf T}$. Since $\mathsf p_1=\varepsilon+\beta+g\ge g$ and $\mathsf p_1=\tilde\varepsilon+\tilde\beta\ge\tilde\varepsilon$, it is at most $\bar{\mathsf b}e^{-g\mathsf T}$ and at most $\bar{\mathsf b}e^{-\tilde\varepsilon\mathsf T}\le\mathsf h_0$.
\end{proof}

\begin{proof}[Proof of \cref{cor:evo-horizon}]
By \cref{thm:evo-record-channel}(a), $\mathsf h\ge\mathsf c\mathsf ue^{-\mathsf T/\tau_{\mathrm{id}}}$, and the right side is at least $\mathsf h_\ast$ exactly when $\mathsf T\le\tau_{\mathrm{id}}\log(\mathsf c\mathsf u/\mathsf h_\ast)$. Without a record, $\mathsf h_0=(\bar{\mathsf b}+(1-\bar{\mathsf b})\mathsf f)e^{-\tilde\varepsilon\mathsf T}\ge\mathsf h_\ast$ is the stated condition on $\mathsf T$ after taking logarithms.
\end{proof}

\begin{proof}[Proof of \cref{prop:evo-variance}]
(a) With the indicators $\xi_i$ of the proof of \cref{prop:evo-transmission}, the deviations of the daughter's birth state from its conditional mean are
\begin{equation*}
\bar\mu'_0-\E[\bar\mu'_0\mid\mathcal H]=\frac1n\sum_i(\xi_i-\bar{\mathsf b})\mu_i+(1-\bar{\mathsf b})\mathsf f\zeta+\frac1n\sum_i(1-\xi_i)\nu_i,\qquad z'_0-\E[z'_0\mid\mathcal H]=\zeta.
\end{equation*}
The founder set, $\zeta$ and the $\nu_i$ are independent of each other and of $\mathcal H$, and $\zeta$ and the $\nu_i$ have zero mean, so the three terms of the first deviation are uncorrelated given $\mathcal H$. For a uniformly drawn set of $\mathsf b$ positions out of $n\ge2$, $\Var\xi_i=\bar{\mathsf b}(1-\bar{\mathsf b})$ and, for $i\ne j$, $\Cov(\xi_i,\xi_j)=\frac{\mathsf b(\mathsf b-1)}{n(n-1)}-\bar{\mathsf b}^2=-\bar{\mathsf b}(1-\bar{\mathsf b})/(n-1)$. Hence
\begin{align*}
\Var\Bigl(\frac1n\sum_i\xi_i\mu_i\Bigm|\mathcal H\Bigr)&=\frac{\bar{\mathsf b}(1-\bar{\mathsf b})}{n^2}\Bigl[\sum_i\mu_i^2-\frac{(\sum_i\mu_i)^2-\sum_i\mu_i^2}{n-1}\Bigr]\\
&=\frac{\bar{\mathsf b}(1-\bar{\mathsf b})}{n^2}\frac{n}{n-1}\sum_i(\mu_i-\bar\mu)^2=\frac{\bar{\mathsf b}(1-\bar{\mathsf b})}{n-1}\mathsf v_{\mathrm w}.
\end{align*}
The third term has variance $\frac1{n^2}\E[\sum_i(1-\xi_i)]\mathsf s_\nu^2=(1-\bar{\mathsf b})\mathsf s_\nu^2/n$ and the second $(1-\bar{\mathsf b})^2\mathsf f^2\mathsf s_z^2$; the covariance of the two deviations is $(1-\bar{\mathsf b})\mathsf f\mathsf s_z^2$, and the variance of the second deviation is $\mathsf s_z^2$. This is \eqref{eq:evo-noise}.

(b) Write $\mathsf x'(\mathsf T)-\mathsf M\mathsf x=\mathsf G(\mathsf T)(\mathsf x'_0-\mathsf F\mathsf x)+(\mathsf x'(\mathsf T)-\mathsf G(\mathsf T)\mathsf x'_0)$. The second term has zero conditional mean given $\mathcal H$ and the daughter's birth state (proof of \cref{prop:evo-transmission}(a)), and the first is a function of these, so the two are uncorrelated given $\mathcal H$, and $\Cov(\mathsf x'(\mathsf T)\mid\mathcal H)=\mathsf G(\mathsf T)\mathsf N_{\mathrm{fis}}\mathsf G(\mathsf T)^\top+\mathsf N_{\mathrm{rel}}$ with $\mathsf N_{\mathrm{rel}}$ the conditional covariance of the second term. Contracting with $\mathsf l$ gives the conditional variance of $y'$. Since $\E[y'\mid\mathcal H]=\mathsf hy$, the law of total variance, applied to a parent drawn from the parental distribution and its one daughter, gives $\Var(y')=\Var(\E[y'\mid\mathcal H])+\E[\Var(y'\mid\mathcal H)]=\mathsf h^2\mathsf V+\mathsf s^2$. Here $\Var$ is taken over the draw of the parent and over the noises of fission and relaxation; for a finite population of $N\ge2$ parents whose daughters' deviations $y'_j-\mathsf hy_j$ are uncorrelated across parents given $\mathcal H$, the sample variance of the daughters' traits with divisor $N-1$ has conditional expectation $\mathsf h^2$ times the parents' sample variance plus $\frac1N\sum_j\Var(y'_j\mid\mathcal H)$, and its realized value fluctuates about that expectation. The statement claims nothing more for a finite population. If $\mathsf s^2$ is the same in every generation, the recursion $\mathsf V_{k+1}=\mathsf h^2\mathsf V_k+\mathsf s^2$ with $\mathsf h<1$ converges to $\mathsf s^2/(1-\mathsf h^2)$. At that value with $\mathsf s^2>0$ one has $\Var(y')=\Var(y)=\mathsf V>0$ and $\Cov(y',y)=\mathsf h\mathsf V$ by \cref{prop:evo-transmission}(a), so $\Corr(y',y)=\mathsf h\mathsf V/(\mathsf V\mathsf V)^{1/2}=\mathsf h$; for $\mathsf s^2=0$ the limit is $0$ and the correlation is undefined.

(c) In the variant without a record, by the proof of \cref{prop:evo-transmission}(d), $\bar\mu'_0-\E[\bar\mu'_0\mid\mathcal H]=(\bar{\mathsf b}+(1-\bar{\mathsf b})\mathsf f)(m_S-\bar\mu)+\frac1n\sum_i(1-\xi_i)\nu_i$. The computation of (a) with $\frac1{\mathsf b}$ in place of $\frac1n$ gives $\Var(m_S\mid\mathcal H)=\frac{\bar{\mathsf b}(1-\bar{\mathsf b})}{\mathsf b^2}\frac{n^2}{n-1}\mathsf v_{\mathrm w}=\frac{\mathsf v_{\mathrm w}}{\mathsf b}\frac{n-\mathsf b}{n-1}$, and the two terms are uncorrelated, which gives the stated variance. Finally $\mathsf h_0=(\mathsf f+(1-\mathsf f)\mathsf b/n)e^{-\tilde\varepsilon\mathsf T}$ is nondecreasing in $\mathsf b$, strictly increasing when $\mathsf f<1$.
\end{proof}

\begin{proposition}[Founders against recruits at long generation times]\label{prop:evo-bottleneck}
Let $\mathsf P(\bar{\mathsf b})=(1-\mathsf u)\bar{\mathsf b}+\mathsf u\mathsf c\mathsf k$ be the limit of $\mathsf he^{\mathsf T/\tau_{\mathrm{id}}}$ in \cref{thm:evo-record-channel}(a), regarded as a function of the founder fraction. It is affine in $\bar{\mathsf b}$ with slope
\begin{equation}\label{eq:evo-bottleneck-slope}
\frac{\mathrm d\mathsf P}{\mathrm d\bar{\mathsf b}}=\frac{\omega(1-\mathsf c\mathsf f)+\delta-\varpi_-}{\varpi_+-\varpi_-}.
\end{equation}
Consequently there is a $\mathsf T_0$ such that for all $\mathsf T\ge\mathsf T_0$ the heritability $\mathsf h$ is strictly increasing in $\mathsf b\in\{0,\dots,n\}$ when $\varpi_-<\delta+\omega(1-\mathsf c\mathsf f)$ and strictly decreasing when $\varpi_->\delta+\omega(1-\mathsf c\mathsf f)$. When $\mathsf c=\mathsf f=1$ the second case holds exactly when $\delta<\tilde\varepsilon$.
\end{proposition}

\begin{proof}
By \eqref{eq:evo-lF}, $\mathsf k$ is affine in $\bar{\mathsf b}$ with slope $-\mathsf f\omega/(\mathsf p_1-\varpi_-)$, so $\mathrm d\mathsf P/\mathrm d\bar{\mathsf b}=(1-\mathsf u)-\mathsf u\mathsf c\mathsf f\omega/(\mathsf p_1-\varpi_-)=[(\mathsf p_2-\varpi_-)-\mathsf c\mathsf f\omega]/(\varpi_+-\varpi_-)$ by \cref{lem:evo-semigroup}(i), and $\mathsf p_2=\omega+\delta$. For each of the finitely many values $\mathsf b\in\{0,\dots,n\}$, $e^{\varpi_-\mathsf T}\mathsf h\to\mathsf P(\mathsf b/n)$ as $\mathsf T\to\infty$ by \cref{thm:evo-record-channel}(a). When the slope is nonzero the limits are strictly ordered in $\mathsf b$, so there is a common $\mathsf T_0$ beyond which the values of $\mathsf h$ are ordered in the same way. When $\mathsf c=\mathsf f=1$ the slope has the sign of $\delta-\varpi_-$, which is negative exactly when $\delta<\tilde\varepsilon$ by \cref{lem:evo-semigroup}(v).
\end{proof}

The mechanism is visible in the slow mode $r$ of \cref{lem:evo-semigroup}(iii). Replacing one recruit by one founder changes $\ell^\top\mathsf F\mathsf x$ by $\omega(\bar\mu-\mathsf f\mathsf cz)/n$, and on the slow mode $\bar\mu/z=\tilde\beta/(\mathsf p_1-\varpi_-)=(\mathsf p_2-\varpi_-)/\omega$. When the record decays more slowly than the members forget, the record coordinate of the slow mode exceeds the members' mean, and a recruit who reads a faithful copy of the record carries more of the heritable trait than a founder does. The comparison of \cref{chk:evo-bottleneck} shows the inversion at short generation times as well in one instance, which the proposition does not cover.

\begin{lemma}[Covariance identity]\label{lem:evo-price-identity}
Let $y_1,\dots,y_N$ be real numbers and $\mathsf W_1,\dots,\mathsf W_N$ nonnegative numbers with $\bar{\mathsf W}=\frac1N\sum_j\mathsf W_j>0$. For each $j$ with $\mathsf W_j>0$ let $\bar y'_j$ be a real number and $\Delta y_j=\bar y'_j-y_j$, and put $\Delta y_j=0$ when $\mathsf W_j=0$. With $\bar y=\frac1N\sum_jy_j$, $\bar y'=\sum_j\mathsf W_j\bar y'_j/\sum_j\mathsf W_j$ and $\Cov_j(\mathsf W_j,y_j)=\frac1N\sum_j\mathsf W_jy_j-\bar{\mathsf W}\bar y$,
\begin{equation}\label{eq:evo-price-identity}
\bar{\mathsf W}(\bar y'-\bar y)=\Cov_j(\mathsf W_j,y_j)+\frac1N\sum_j\mathsf W_j\Delta y_j.
\end{equation}
\end{lemma}

\begin{proof}
$N\bar{\mathsf W}\bar y'=\sum_j\mathsf W_j\bar y'_j=\sum_j\mathsf W_jy_j+\sum_j\mathsf W_j\Delta y_j$. Dividing by $N$ and subtracting $\bar{\mathsf W}\bar y$ gives \eqref{eq:evo-price-identity}.
\end{proof}

\begin{proof}[Proof of \cref{prop:evo-price}]
Apply \cref{lem:evo-price-identity} with $y_j=\mathsf a^\top\mathsf x_j$ and $\bar y'_j$ the mean of $\mathsf a^\top\mathsf x'(\mathsf T)$ over the $\mathsf W_j$ daughters of parent $j$, and divide by $\bar{\mathsf W}$: $\bar y'-\bar y=\Cov_j(\mathsf w_j,y_j)+\frac1N\sum_j\mathsf w_j\Delta y_j$. The numbers $\mathsf W_j$ and $y_j$ are $\mathcal H$-measurable, and each daughter of parent $j$ is produced by a fission of \cref{def:evo-fission}, so by \cref{prop:evo-transmission}(a) $\E[\bar y'_j\mid\mathcal H]=\mathsf a^\top\mathsf M\mathsf x_j$ and $\E[\Delta y_j\mid\mathcal H]=\mathsf a^\top(\mathsf M-I)\mathsf x_j$ when $\mathsf W_j>0$. Since $\frac1N\sum_j\mathsf w_j=1$, $\frac1N\sum_j\mathsf w_j\mathsf a^\top(\mathsf M-I)\mathsf x_j=\Cov_j(\mathsf w_j,\mathsf a^\top(\mathsf M-I)\mathsf x_j)+\mathsf a^\top(\mathsf M-I)\bar{\mathsf x}$, and adding $\Cov_j(\mathsf w_j,\mathsf a^\top\mathsf x_j)$ gives \eqref{eq:evo-price-general}. For $\mathsf a=\mathsf l$, $\mathsf l^\top\mathsf M=\mathsf h\mathsf l^\top$, so the transmission term is $\frac1N\sum_j\mathsf w_j(\mathsf h-1)y_j=-(1-\mathsf h)(\bar y+\mathsf d_{\mathrm{sel}})$, and the sum of the two terms is $\mathsf h\mathsf d_{\mathrm{sel}}-(1-\mathsf h)\bar y$, which is positive exactly when $\mathsf h\mathsf d_{\mathrm{sel}}>(1-\mathsf h)\bar y$. If all $\mathsf W_j$ are equal then $\mathsf d_{\mathrm{sel}}=0$ and $\E[\bar y'\mid\mathcal H]=\mathsf h\bar y$; if this holds in every generation, the tower property gives $\E[\bar y_k]=\mathsf h^k\E[\bar y_0]$.
\end{proof}

\begin{proof}[Proof of \cref{prop:evo-balance}]
Write $\mathcal N(\cdot;m,\mathsf V)$ for the normal density.

\emph{Recursion.} For $\mathsf V>0$ the trait density among the parents of the daughters is proportional to $\mathsf W(y)\mathcal N(y;\bar y,\mathsf V)$, the exponential of a quadratic with second-order coefficient $-\frac12(1/\mathsf V+1/\mathsf v_{\mathrm s})$. It is therefore normal with variance $\mathsf V\mathsf v_{\mathrm s}/(\mathsf V+\mathsf v_{\mathrm s})$ and mean $(\bar y/\mathsf V+\mathsf y_\ast/\mathsf v_{\mathrm s})\mathsf V\mathsf v_{\mathrm s}/(\mathsf V+\mathsf v_{\mathrm s})=\bar y+\frac{\mathsf V}{\mathsf V+\mathsf v_{\mathrm s}}(\mathsf y_\ast-\bar y)$; for $\mathsf V=0$ the weighting leaves the point mass at $\bar y$ unchanged and the same formulas hold. A daughter's trait is $\mathsf hy$ plus an independent normal variable of variance $\mathsf s^2$, so the daughters' traits are normal with the mean and variance of \eqref{eq:evo-balance-recursion}. The selection differential is $\mathsf d_{\mathrm{sel}}=\frac{\mathsf V}{\mathsf V+\mathsf v_{\mathrm s}}(\mathsf y_\ast-\bar y)$, and the recursion for the mean is $\bar y'=\mathsf h(\bar y+\mathsf d_{\mathrm{sel}})$, in agreement with \eqref{eq:evo-price}.

\emph{Variance.} Let $\Phi(\mathsf V)=\mathsf h^2\mathsf v_{\mathrm s}\mathsf V/(\mathsf V+\mathsf v_{\mathrm s})+\mathsf s^2$ on $[0,\infty)$. Then
\begin{equation}\label{eq:evo-balance-q}
\Phi(\mathsf V)-\mathsf V=-\frac{\mathsf q(\mathsf V)}{\mathsf V+\mathsf v_{\mathrm s}},\qquad\mathsf q(\mathsf V)=\mathsf V^2+\mathsf V(\mathsf v_{\mathrm s}(1-\mathsf h^2)-\mathsf s^2)-\mathsf s^2\mathsf v_{\mathrm s}.
\end{equation}
The quadratic $\mathsf q$ has roots of product $-\mathsf s^2\mathsf v_{\mathrm s}<0$, hence exactly one positive root, which is $\hat{\mathsf V}$ of \eqref{eq:evo-balance-fixed}; $\mathsf q<0$ on $[0,\hat{\mathsf V})$ and $\mathsf q>0$ on $(\hat{\mathsf V},\infty)$. Since $\Phi\ge\mathsf s^2$, $\mathsf V_k\ge\mathsf s^2$ for $k\ge1$ and $\hat{\mathsf V}=\Phi(\hat{\mathsf V})\ge\mathsf s^2$. On $[\mathsf s^2,\infty)$, $0\le\Phi'(\mathsf V)=\mathsf h^2\mathsf v_{\mathrm s}^2/(\mathsf V+\mathsf v_{\mathrm s})^2\le\theta^2$ with $\theta=\mathsf v_{\mathrm s}/(\mathsf s^2+\mathsf v_{\mathrm s})<1$, so $\lvert\mathsf V_{k+1}-\hat{\mathsf V}\rvert\le\theta^2\lvert\mathsf V_k-\hat{\mathsf V}\rvert$ for $k\ge1$ and $\mathsf V_k\to\hat{\mathsf V}$.

\emph{Mean.} Put $\mathsf i_k=\mathsf V_k/(\mathsf V_k+\mathsf v_{\mathrm s})$, so that $\mathsf i_k\to\mathsf i=\hat{\mathsf V}/(\hat{\mathsf V}+\mathsf v_{\mathrm s})>0$ and $\bar y_{k+1}=\mathsf h(1-\mathsf i_k)\bar y_k+\mathsf h\mathsf i_k\mathsf y_\ast$. The number $1-\mathsf h+\mathsf h\mathsf i$ is positive, because $1-\mathsf h>0$ or $\mathsf h=1$ and $\mathsf i>0$, and $\hat y$ of \eqref{eq:evo-balance-fixed} is the unique solution of $\hat y=\mathsf h(1-\mathsf i)\hat y+\mathsf h\mathsf i\mathsf y_\ast$. Subtracting, $e_k=\bar y_k-\hat y$ satisfies $e_{k+1}=\mathsf h(1-\mathsf i_k)e_k+\mathsf h(\mathsf i_k-\mathsf i)(\mathsf y_\ast-\hat y)$. For $k\ge1$, $\mathsf i_k\ge\mathsf s^2/(\mathsf s^2+\mathsf v_{\mathrm s})$, so $0\le\mathsf h(1-\mathsf i_k)\le\theta$ and $\lvert e_{k+1}\rvert\le\theta\lvert e_k\rvert+\eta_k$ with $\eta_k=\lvert\mathsf i_k-\mathsf i\rvert\lvert\mathsf y_\ast-\hat y\rvert\to0$. The sequence $\lvert e_k\rvert$ is bounded, and its upper limit $L$ satisfies $L\le\theta L$, so $L=0$.

\emph{Dependence on $\mathsf h$ and $\mathsf s^2$.} For $\mathsf V>0$, $\Phi(\mathsf V)$ is strictly increasing in $\mathsf h\in[0,1]$ and in $\mathsf s^2$. If $\Phi_1\le\Phi_2$ are the maps for two parameter values, with strict inequality on $(0,\infty)$, and $\hat{\mathsf V}_1$, $\hat{\mathsf V}_2$ their fixed points, then $\Phi_2(\hat{\mathsf V}_1)-\hat{\mathsf V}_1>\Phi_1(\hat{\mathsf V}_1)-\hat{\mathsf V}_1=0$, and by \eqref{eq:evo-balance-q} the function $\Phi_2(\mathsf V)-\mathsf V$ is positive only below $\hat{\mathsf V}_2$, so $\hat{\mathsf V}_1<\hat{\mathsf V}_2$. Hence $\hat{\mathsf V}$ and $\mathsf i$ are strictly increasing in $\mathsf h$ and in $\mathsf s^2$; they are continuous in both by \eqref{eq:evo-balance-fixed}, and $\hat{\mathsf V}=\mathsf s^2$ at $\mathsf h=0$, where $\mathsf q(\mathsf V)=(\mathsf V-\mathsf s^2)(\mathsf V+\mathsf v_{\mathrm s})$. The retained fraction is $\mathsf R=\hat y/\mathsf y_\ast=\mathsf h\mathsf i/(1-\mathsf h+\mathsf h\mathsf i)$, continuous on $[0,1]$, with $\mathsf R=0$ at $\mathsf h=0$, $\mathsf R=1$ at $\mathsf h=1$ and $\mathsf R>0$ for $\mathsf h>0$. For $0<\mathsf h\le1$, $\mathsf R=[1+(1-\mathsf h)/(\mathsf h\mathsf i)]^{-1}$, and $(1-\mathsf h)/\mathsf h$ is nonnegative and strictly decreasing while $1/\mathsf i$ is positive and strictly decreasing, so their product is strictly decreasing on $(0,1]$ and $\mathsf R$ is strictly increasing there, hence on $[0,1]$. A threshold would be a value $\mathsf h_{\mathrm c}>0$ with $\mathsf R=0$ for $\mathsf h\le\mathsf h_{\mathrm c}$, or a discontinuity of $\mathsf R$; neither exists. For fixed $0<\mathsf h<1$, $\mathsf R$ is strictly increasing in $\mathsf i$, hence in $\mathsf s^2$. From $\hat{\mathsf V}=\Phi(\hat{\mathsf V})\le\mathsf h^2\hat{\mathsf V}+\mathsf s^2$ one gets $\hat{\mathsf V}\le\mathsf s^2/(1-\mathsf h^2)$, so $\hat{\mathsf V}\to0$, $\mathsf i\to0$ and $\mathsf R\le\mathsf h\mathsf i/(1-\mathsf h)\to0$ as $\mathsf s^2\to0$. Finally, for $\mathsf h>0$, $\mathsf R\ge\frac12$ exactly when $(1-\mathsf h)/(\mathsf h\mathsf i)\le1$, that is, when $\mathsf h\mathsf i\ge1-\mathsf h$; at $\mathsf h=0$ both $\mathsf R\ge\frac12$ and $\mathsf h\mathsf i\ge1-\mathsf h$ fail.
\end{proof}

The price of the variance that sustains the response is a load in mean fitness. At the fixed point the mean fitness is
\begin{equation}\label{eq:evo-balance-fitness}
\int\mathsf W(y)\mathcal N(y;\hat y,\hat{\mathsf V})\mathrm dy=\Bigl(\frac{\mathsf v_{\mathrm s}}{\mathsf v_{\mathrm s}+\hat{\mathsf V}}\Bigr)^{1/2}\exp\Bigl(-\frac{(\hat y-\mathsf y_\ast)^2}{2(\mathsf v_{\mathrm s}+\hat{\mathsf V})}\Bigr),
\end{equation}
because $\mathsf W(y)=(2\pi\mathsf v_{\mathrm s})^{1/2}\mathcal N(\mathsf y_\ast;y,\mathsf v_{\mathrm s})$, so the integral is $(2\pi\mathsf v_{\mathrm s})^{1/2}$ times the density at $\mathsf y_\ast$ of the sum of independent normal variables with laws $\mathcal N(\hat y,\hat{\mathsf V})$ and $\mathcal N(0,\mathsf v_{\mathrm s})$. The two bounds $\mathsf s^2\le\hat{\mathsf V}\le\mathsf s^2/(1-\mathsf h^2)$ of the proof, the second for $\mathsf h<1$, also give the weak-selection limit: as $\mathsf v_{\mathrm s}\to\infty$ with $\mathsf h<1$ fixed, $\mathsf v_{\mathrm s}/(\hat{\mathsf V}+\mathsf v_{\mathrm s})\to1$ and $\hat{\mathsf V}\to\mathsf s^2/(1-\mathsf h^2)$, the neutral value of \cref{prop:evo-variance}(b). At $\mathsf h=1$ the fixed point is $\hat{\mathsf V}=\frac12[\mathsf s^2+(\mathsf s^4+4\mathsf s^2\mathsf v_{\mathrm s})^{1/2}]$ and $\hat y=\mathsf y_\ast$.

\begin{lemma}[Monotonicity in the entries]\label{lem:evo-monotone}
If $\mathsf Q$ and $\mathsf Q'$ are $2\times2$ matrices with $0\le\mathsf Q\le\mathsf Q'$, then $\rho(\mathsf Q)\le\rho(\mathsf Q')$. Consequently $\mathsf h$ is nondecreasing in the copy fidelity $\mathsf c$, and under \cref{ass:evo-gaussian} with $\mathsf s^2$ held fixed the retained fraction is nondecreasing in $\mathsf c$.
\end{lemma}

\begin{proof}
Let $\chi_{\mathsf Q}(\lambda)=(\lambda-q_{11})(\lambda-q_{22})-q_{12}q_{21}$ and $\lambda'=\rho(\mathsf Q')$. By \cref{lem:evo-perron}(i), $\lambda'\ge\max\{q'_{11},q'_{22}\}\ge\max\{q_{11},q_{22}\}$, so $0\le\lambda'-q'_{ii}\le\lambda'-q_{ii}$, and $q'_{12}q'_{21}\ge q_{12}q_{21}$; hence $\chi_{\mathsf Q}(\lambda')\ge\chi_{\mathsf Q'}(\lambda')=0$. The parabola $\chi_{\mathsf Q}$ is negative strictly between its roots and its larger root $\rho(\mathsf Q)$ is at least its vertex $\frac12(q_{11}+q_{22})$; since $\lambda'$ is at least the vertex and $\chi_{\mathsf Q}(\lambda')\ge0$, $\lambda'\ge\rho(\mathsf Q)$. By \eqref{eq:evo-M-entries} the entries $\mathsf M_{12}$ and $\mathsf M_{22}$ are proportional to $\mathsf c$ and the others do not depend on it, which gives the monotonicity of $\mathsf h$; the last claim follows from \cref{prop:evo-balance}.
\end{proof}

The founder fraction has no such monotonicity: $\mathsf M_{11}$ and $\mathsf M_{21}$ increase with $\bar{\mathsf b}$ while $\mathsf M_{12}$ and $\mathsf M_{22}$ decrease when $\mathsf f\mathsf c>0$, and \cref{prop:evo-bottleneck} shows that either direction occurs.

\begin{proposition}[Directional selection]\label{prop:evo-directional}
In \cref{ass:evo-gaussian} replace the fitness by $\mathsf W(y)=e^{\mathsf ry}$ with $\mathsf r\in\R$. Then $\bar y'=\mathsf h(\bar y+\mathsf r\mathsf V)$ and $\mathsf V'=\mathsf h^2\mathsf V+\mathsf s^2$. For $\mathsf h<1$ the recursion converges from every initial condition to
\begin{equation}\label{eq:evo-directional}
\hat{\mathsf V}=\frac{\mathsf s^2}{1-\mathsf h^2},\qquad\hat y=\frac{\mathsf h\mathsf r\mathsf s^2}{(1-\mathsf h)^2(1+\mathsf h)}.
\end{equation}
For $\mathsf h=1$ and $\mathsf r\ne0$ the variance and the absolute value of the mean grow without bound.
\end{proposition}

\begin{proof}
The density proportional to $e^{\mathsf ry}\mathcal N(y;\bar y,\mathsf V)$ is $\mathcal N(y;\bar y+\mathsf r\mathsf V,\mathsf V)$, by completing the square, and transmission multiplies the mean by $\mathsf h$ and maps the variance to $\mathsf h^2\mathsf V+\mathsf s^2$. For $\mathsf h<1$ the variance recursion is an affine contraction with fixed point $\mathsf s^2/(1-\mathsf h^2)$, and $\bar y_{k+1}=\mathsf h\bar y_k+\mathsf h\mathsf r\mathsf V_k$ is an affine contraction with convergent forcing, so by the argument used for the mean in the proof of \cref{prop:evo-balance} it converges to $\mathsf h\mathsf r\hat{\mathsf V}/(1-\mathsf h)$. For $\mathsf h=1$, $\mathsf V_k=\mathsf V_0+k\mathsf s^2$ and $\bar y_k=\bar y_0+\mathsf r\sum_{j<k}\mathsf V_j$.
\end{proof}

With a fitness that is unbounded in the trait the stationary displacement again has the sign of $\mathsf r$ for every $\mathsf h>0$ and diverges as $\mathsf h\to1$; a threshold is absent here too.

\begin{proposition}[A two-type model with a threshold]\label{prop:evo-two-type}
Let an infinite population of collectives consist of carriers of a convention and blanks. A carrier founds $\mathsf m>1$ daughters for each daughter of a blank, a daughter of a carrier is a carrier with probability $\mathsf q\in(0,1]$ and a blank otherwise, and every daughter of a blank is a blank. The frequency $p$ of carriers obeys
\begin{equation}\label{eq:evo-two-type}
p'=\frac{\mathsf m\mathsf qp}{1+(\mathsf m-1)p}.
\end{equation}
If $\mathsf m\mathsf q\le1$ then $p_k\to0$ from every $p_0\in[0,1]$. If $\mathsf m\mathsf q>1$ and $p_0>0$ then $p_k\to\hat p=(\mathsf m\mathsf q-1)/(\mathsf m-1)$.
\end{proposition}

\begin{proof}
The daughters of carriers are a fraction $\mathsf mp/(\mathsf mp+1-p)$ of all daughters, and a fraction $\mathsf q$ of them are carriers, which is \eqref{eq:evo-two-type}. The map $\Psi(p)=\mathsf m\mathsf qp/(1+(\mathsf m-1)p)$ is increasing on $[0,1]$ with values in $[0,\mathsf q]$, and $\Psi(p)/p=\mathsf m\mathsf q/(1+(\mathsf m-1)p)$ is strictly decreasing in $p$. Its fixed points are $0$ and the solution $\hat p$ of $1+(\mathsf m-1)p=\mathsf m\mathsf q$, which lies in $(0,1]$ exactly when $\mathsf m\mathsf q>1$. If $\mathsf m\mathsf q\le1$, then $\Psi(p)<p$ for every $p>0$, so $p_k$ decreases to a fixed point in $[0,p_0]$, which must be $0$. If $\mathsf m\mathsf q>1$ and $0<p<\hat p$, then $\Psi(p)/p>\Psi(\hat p)/\hat p=1$ and $\Psi(p)<\Psi(\hat p)=\hat p$, so $p_k$ increases to a positive fixed point, which is $\hat p$; if $p>\hat p$ the inequalities are reversed and $p_k$ decreases to $\hat p$.
\end{proof}

This model lies outside the linear-Gaussian setting and is included for contrast with \cref{prop:evo-balance}. The threshold $\mathsf m\mathsf q=1$ exists because the advantage of a carrier is bounded and the blank state is absorbing. In the Gaussian model the loss $(1-\mathsf h)y$ and the selection differential are both linear in the displacement and the transmission noise restores variance in every generation, so selection always retains a positive fraction. A retention probability $\mathsf q$ would have to be derived from a bistable convention dynamics of the kind conjectured in \cref{conj:agy-conformity-identity}, which is \cref{open:evo-threshold}.

\begin{example}[A durable record with naive recruits]\label{ex:evo-worked}
Take the residence time as the unit of time, $g=1$, and let $\varepsilon=0.01$, $\beta=2$, $\mathsf f=0$, $\omega=0.05$, $\delta=0$, $\mathsf c=0.95$, $n=20$ and $\mathsf b=2$. Then $\tilde\varepsilon=1.01$, $\tilde\beta=2$, $\tau_{\mathrm{id}}=60.27$ against $1/\tilde\varepsilon=0.990$ without a record, $\mathsf u=0.9890$, $\mathsf k=1$ and the limit of \cref{thm:evo-record-channel}(a) is $(1-\mathsf u)\bar{\mathsf b}+\mathsf u\mathsf c=0.9406$. \Cref{tab:evo-worked} lists the heritability at four generation times, together with the value $\mathsf h_0$ without a record, the value $\bar{\mathsf b}\mathsf G_{11}(\mathsf T)$ when the record is blank at birth, and the coefficient $\bar{\mathsf b}e^{-\mathsf p_1\mathsf T}$ of \cref{thm:evo-record-channel}(e) for a record written by non-members. By \cref{cor:evo-horizon}, $\mathsf h\ge\frac12$ for every generation time up to $\tau_{\mathrm{id}}\log(2\mathsf c\mathsf u)=38.0$ residence times, whereas $\mathsf h_0\le\bar{\mathsf b}=0.1$ at every generation time. Naive recruits and two founders out of twenty leave almost nothing to inherit through the members, and the copied record carries the convention across the fission.
\end{example}

\begin{table}[t]
\centering\small
\begin{tabular}{@{}rccccc@{}}
\toprule
$\mathsf T$ & $\mathsf h$ & $\mathsf he^{\mathsf T/\tau_{\mathrm{id}}}$ & $\mathsf h_0$ (no record) & blank record ($\mathsf c=0$) & external record\\
\midrule
$1$ & $0.9255$ & $0.9410$ & $3.6\times10^{-2}$ & $5.8\times10^{-3}$ & $4.9\times10^{-3}$\\
$5$ & $0.8657$ & $0.9406$ & $6.4\times10^{-4}$ & $1.0\times10^{-3}$ & $2.9\times10^{-8}$\\
$20$ & $0.6750$ & $0.9406$ & $1.7\times10^{-10}$ & $7.9\times10^{-4}$ & $7.2\times10^{-28}$\\
$60$ & $0.3476$ & $0.9406$ & $4.8\times10^{-28}$ & $4.1\times10^{-4}$ & $3.7\times10^{-80}$\\
\bottomrule
\end{tabular}
\caption{Heritability of the convention in \cref{ex:evo-worked} as a function of the generation time $\mathsf T$, in units of the residence time. The third column converges to $0.9406$; the last three columns are the transmission coefficients without a record, with a record that is blank at birth, and with a record written only by non-members.}
\label{tab:evo-worked}
\end{table}

\begin{remark}[Comparison with tradeoff breaking]\label{rem:evo-bourrat-detail}
\citet{bourrat2022tradeoff} examine the description of an evolutionary transition in individuality as a decoupling of the fitness of collectives from the fitness of their particles. Their result is that when both fitnesses are measured over commensurable time intervals, in a population whose distribution of collective sizes has stabilized, the two coincide, so the fitnesses cannot decouple. They propose to describe transitions through traits and the tradeoffs among them, with the breaking of a tradeoff as a marker of a transition when it is observed in the context of one; they state that tradeoff breaking can occur outside transitions and that such an event taken alone does not identify a transition. Three points of contact with the fission model are exact. First, the model has no particle-level fitness: members are recruited and replaced and are not born of members, so the comparison of the two fitnesses cannot be made inside it; if one counts occupied positions instead, there are $n$ of them in every collective at every census, their number grows at the rate of the number of collectives, and the agreement of the two rates holds trivially and tests nothing. Second, the construction is trait-based in the sense those authors recommend: the fitness $\mathsf W$ is a declared function of a collective trait, and \cref{prop:evo-price,prop:evo-balance} follow the trait. Third, the model has a single trait, so it contains no tradeoff and cannot express the breaking of one. The following suggestion is interpretive and is not a result. The manuscript does contain two collective quantities governed by the same parameter, the read rate $\beta$, one of which can improve while the other cannot. Differentiating $\chi(\varpi_-)=0$ in \eqref{eq:evo-charpoly} with respect to $\beta$, with $\partial\chi/\partial\beta=\delta-s$ and $\chi'(\varpi_-)=-(\varpi_+-\varpi_-)$, gives $\partial\varpi_-/\partial\beta=(\delta-\varpi_-)/(\varpi_+-\varpi_-)$, which is negative exactly when $\delta<\tilde\varepsilon$ by \cref{lem:evo-semigroup}(v). When the record decays more slowly than the members forget, reading a member-written record more strongly therefore lengthens the identity time of a convention (compare \cref{cor:hie-identity}) and lowers the exponential rate at which its heritability decays with the generation time in \eqref{eq:evo-sandwich}. The heritability at a fixed generation time need not increase with $\beta$: for a record that is blank at birth, $\mathsf h=\bar{\mathsf b}\mathsf G_{11}(\mathsf T)=\bar{\mathsf b}(1-\mathsf p_1\mathsf T)+O(\mathsf T^2)$ decreases with $\beta$ at short generation times (\cref{chk:evo-heritability}). For a grounded coordinate the common-mode factor that bounds every member's stationary error from below is nondecreasing in $\beta$ (\cref{thm:hie-floor}(c), where $\beta$ is the per-round read gain of the discrete-time record of \cref{lem:hie-record-blind} and $\beta,\omega>0$), so stronger reading never lowers that floor. Whether fitness functions of these two quantities define a tradeoff in the sense of those authors, and whether a change of organization can break it, is \cref{open:evo-tradeoff}.
\end{remark}
